\documentclass[11pt]{article}

\usepackage[a4paper,margin=30mm]{geometry}
\usepackage{amsmath,amsthm,amsfonts,mathtools}

% Use one optical size of Latin Modern Roman for XeTeX text at every size.
\ifdefined\XeTeXversion
  \DeclareFontShape{TU}{lmr}{m}{n}
    {<->"[lmroman10-regular]:mapping=tex-text;"}{}
  \DeclareFontShape{TU}{lmr}{m}{it}
    {<->"[lmroman10-italic]:mapping=tex-text;"}{}
\fi

\newcommand{\E}{\mathbb{E}}
\newcommand{\Pp}{\mathbf{P}}
\newcommand{\R}{\mathbf{R}}
\newcommand{\eqdef}{\mathrel{\vcentcolon=}}
\newcommand{\arginf}{\mathop{\mathrm{arg\,inf}}}

% Notation for the population and empirical objectives.
\newcommand{\rhohu}{\rho}
\newcommand{\rhohum}{\phi}
\newcommand{\Kern}{K}
\newcommand{\Signal}{X}
\newcommand{\BiasedSignal}{X^{\mathrm b}}
\newcommand{\Data}{Y}
\newcommand{\IMFs}{S}
\newcommand{\IMFt}{\widetilde S}
\newcommand{\IF}{F}
\newcommand{\IFP}{F_{\circ}}
\newcommand{\DFN}{D}
\newcommand{\fs}{f}
\newcommand{\fp}{f_{\circ}}
\newcommand{\xs}{x^{*}}
\newcommand{\xp}{x_{\circ}}
\newcommand{\xx}{x}

\newtheorem{proposition}{Proposition}[section]
\newtheorem{lemma}[proposition]{Lemma}
\newtheorem{corollary}[proposition]{Corollary}
\newtheorem*{remark}{Remark}

\begin{document}
\begin{center}
  {\Large The first smooth-absolute IRMF component\par}
  \vspace{5pt}
  {\large Sampling error, target shift, and biased signal error\par}
\end{center}

\setcounter{section}{3}
\section{Theoretical guarantees}

The first IRMF component fits a local constant to noisy observations using
Gaussian-smoothed absolute loss.  We first bound the sampling error around
the population target.  We then bound the deterministic shift from that
target to the biased signal, the output of the same filter on noiseless
data.  Adding these two terms gives the error relative to the biased signal.

\subsection*{The three fits and what their differences measure}
Let \(Y_u=X_u+\varepsilon_u\), where \(X_u\) is the clean signal and
\(\varepsilon_u\) is observation noise.  At a location \(t\), use fixed
weights \(w_{u,t}\geq0\) with positive total weight.  All three fits use
these same weights and the same loss
\[
  \rho_H(r)=\E_Z|r+HZ|,\qquad H>0,\qquad Z\sim N(0,1),
\]
where the auxiliary variable \(Z\) is independent of the observation noise.
Define
\begin{align*}
  \widetilde S_t
    &\eqdef\arginf_{z\in\R}\sum_u w_{u,t}\rho_H(Y_u-z),
    &&\text{observed component},\\
  S_t
    &\eqdef\arginf_{z\in\R}\sum_u w_{u,t}
      \E_\varepsilon\rho_H(X_u+\varepsilon_u-z),
    &&\text{population target},\\
  \BiasedSignal_t
    &\eqdef\arginf_{z\in\R}\sum_u w_{u,t}\rho_H(X_u-z),
    &&\text{biased signal}.
\end{align*}
The observed component \(\widetilde S_t\) is what the algorithm computes
from the noisy data.  The biased signal \(\BiasedSignal_t\) is what the same algorithm
would compute if the clean signal were available.  The population target
\(S_t\) minimizes the noisy loss after that loss has been averaged over
possible noise realizations.

Only \(\widetilde S_t\) changes with the realized noise.  For a fixed signal,
design, and noise law, \(S_t\) and \(\BiasedSignal_t\) are deterministic.  Averaging a
loss and then minimizing it is generally different from averaging its
minimizer, so \(S_t\) need not equal \(\E\widetilde S_t\).

\begin{description}
\item[\(\widetilde S_t-S_t\): sampling error.]
The displacement of the observed fit from its population target.  This
finite-sample error need not have mean zero.
\item[\(S_t-\BiasedSignal_t\): target shift.]
The change in the minimizer caused by averaging over observation noise.
It can be nonzero even for mean-zero noise because the loss is nonlinear.
\item[\(\widetilde S_t-\BiasedSignal_t\): biased signal error.]
The total change in the extracted component caused by adding observation
noise to the clean signal.
\end{description}
These differences satisfy the exact identity
\[
  \boxed{\widetilde S_t-\BiasedSignal_t
    =\underbrace{\widetilde S_t-S_t}_{\text{sampling error}}
     +\underbrace{S_t-\BiasedSignal_t}_{\text{target shift}}.}
\]
\newpage
We call \(\BiasedSignal_t\) the \emph{biased signal} because local filtering
can change the signal value even without observation noise.  The difference
\(\BiasedSignal_t-X_t\) is deterministic filtering distortion; the name does
not identify \(\BiasedSignal_t\) with the expectation of an estimator.
The biased signal error measures deviation from \(\BiasedSignal_t\), not
from \(X_t\).  Recovering \(X_t\) would also require controlling
\(\BiasedSignal_t-X_t\).

\subsection*{From sampling error to biased signal error}
Put \(\Delta_t=S_t-\BiasedSignal_t\).  The argument has three steps:
\begin{enumerate}
\item Bound the sampling error using the expansion
\[
  \widetilde S_t-S_t=F_t^{-1}\nabla_t+R_t,
  \qquad |\widetilde S_t-S_t|\leq\mathcal E_x.
\]
Here \(F_t\) is the population curvature at \(S_t\), \(\nabla_t\) is the
centered empirical score there, and \(R_t\) is the remainder.
Proposition~\ref{prop:main} gives the general expansion;
Corollary~\ref{cor:gaussian} specifies \(\mathcal E_x\) under Gaussian noise.
\item Bound the target shift, which contains no sampling fluctuation:
\[
  |\Delta_t|=|S_t-\BiasedSignal_t|\leq\beta_t.
\]
Proposition~\ref{prop:shift} obtains this deterministic bound by comparing
the score equations for the two deterministic targets.
\item Add the shift to obtain the biased signal error:
\[
  \widetilde S_t-\BiasedSignal_t
    =F_t^{-1}\nabla_t+R_t+\Delta_t,
  \qquad |\widetilde S_t-\BiasedSignal_t|\leq\mathcal E_x+\beta_t.
\]
Corollary~\ref{cor:biased} combines the first two steps on the same
probability event.  Subtracting the oracle term \(F_t^{-1}\nabla_t\)
leaves \(R_t+\Delta_t\), so it does not remove the target shift.
\end{enumerate}

For orientation, let \(n\) be the number of observations, \(m\) the number
of evaluation points, and \(h\) the spatial window radius.  Set
\(z_x=x+\log(4m)\), where \(x\geq0\) sets the failure probability \(e^{-x}\).
For independent Gaussian noise, fixed loss and noise scales, regular windows,
and normalized population curvature bounded away from zero,
\[
  \mathcal E_x=O\!\left(\sqrt{z_x/(nh)}\right),\qquad
  \mathcal R_x=O\!\left(z_x/(nh)\right),
\]
where \(\mathcal R_x\) bounds \(\max_t|R_t|\).
Corollary~\ref{cor:sample-size} gives sufficient local sample-size conditions.
The shift \(\Delta_t\) can persist as the sampling error decreases.
All proofs are in Appendix~\ref{sec:proofs}.

\clearpage
\subsection{Sampling error: model and expansion}\label{sec:model}

Let \(u\) range over \(n\) deterministic observation locations in \([0,1]\).
Let \(\mathcal T_n\subset[0,1]\) be a nonempty finite set of evaluation points
and put \(m=|\mathcal T_n|\).  The signal, design, bandwidth, and weights are
deterministic.  For \(t\in\mathcal T_n\), write
\begin{equation}
  w_{u,t}\eqdef \Kern_h(t-u),\qquad
  W_t\eqdef\sum_u w_{u,t},\qquad
  Q_t^2\eqdef\sum_u w_{u,t}^2,\qquad
  k_t\eqdef\max_u w_{u,t}.
  \label{eq:weights}
\end{equation}
Assume \(w_{u,t}\geq0\) and \(W_t>0\).  Suppose that
\(\Data_u=\Signal_u+\varepsilon_u\), where the errors are independent and
identically distributed with \(\E|\varepsilon_u|<\infty\), and let
\begin{equation}
  \rhohum(z)\eqdef\E\rhohu(z+\varepsilon).
  \label{eq:phi}
\end{equation}
The argument below also allows independent, non-identically distributed errors
with finite first moments if \(\rhohum(\Signal_u-x)\) in each summand is replaced by
\(\E\rhohu(\Data_u-x)\).

Assume that \(\rhohu\in C^3(\R)\) is convex and coercive, meaning
\(\rhohu(z)\to\infty\) as \(|z|\to\infty\), and, for finite constants
\(L_1,L_2,L_3\),
\begin{equation}
  |\rhohu'(z)|\leq L_1,\qquad
  0\leq\rhohu''(z)\leq L_2,\qquad
  |\rhohu'''(z)|\leq L_3,
  \qquad z\in\R.
  \label{eq:derivative-bounds}
\end{equation}
For the Gaussian-smoothed absolute-value contrast, with a deterministic \(H>0\),
\begin{align}
  \rho_H(r)&\eqdef\E_Z|r+HZ|,\qquad Z\sim N(0,1),\notag\\
  \rho_H'(r)&=\operatorname{erf}\!\left(\frac{r}{\sqrt2H}\right),\qquad
  \rho_H''(r)=\frac{\sqrt{2/\pi}}{H}e^{-r^2/(2H^2)},\notag\\
  \rho_H'''(r)&=-\frac{\sqrt{2/\pi}\,r}{H^3}e^{-r^2/(2H^2)}.
  \label{eq:smooth-abs}
\end{align}
One may take the conservative constants
\begin{equation}
  L_1=1,\qquad
  L_2=\frac{2}{H\sqrt{2\pi}},\qquad
  L_3=\frac{2}{H^2\sqrt{2\pi}}.
  \label{eq:rhoH-constants}
\end{equation}

The Lipschitz bound \(|\rho'|\leq L_1\) and the finite first noise moment make
the expected objectives finite.  The inequality
\(\E\rho(Y_u-z)\geq\rho(-z)-L_1\E|Y_u|\) also makes each population
objective coercive.  Thus both empirical and population minimizers exist.
Define the empirical fit \(\widetilde S_t\), its population target
\(S_t\), the population curvature \(F_t\), and the centered score
\(\nabla_t\) by
\begin{align}
  \IMFt_t
  &\eqdef \arginf_{z\in\R}\sum_u\rhohu(\Data_u-z)w_{u,t},
  &
  \IMFs_t
  &\eqdef \arginf_{z\in\R}\sum_u\E\rhohu(\Data_u-z)w_{u,t},
  \label{eq:estimators}\\
  \IF_t
  &\eqdef\sum_u\rhohum''(\Signal_u-\IMFs_t)w_{u,t},
  &
  \DFN_t
  &\eqdef\IF_t^{1/2},
  \label{eq:curvature}\\
  \nabla_t
  &\eqdef\sum_u\left\{
       \rhohu'(\Data_u-\IMFs_t)
       -\E\rhohu'(\Data_u-\IMFs_t)
     \right\}w_{u,t}.
  \label{eq:nabla}
\end{align}
Also define the biased signal and target shift by
\begin{equation}
  \BiasedSignal_t\eqdef\arginf_{z\in\R}\sum_u\rho(X_u-z)w_{u,t},
  \qquad\Delta_t\eqdef S_t-\BiasedSignal_t.
  \label{eq:biased-signal}
\end{equation}
Assume this minimizer is unique; this is automatic for \(\rho_H\).
The two signed errors satisfy the exact identity
\begin{equation}
  \widetilde S_t-\BiasedSignal_t=(\widetilde S_t-S_t)+\Delta_t.
  \label{eq:target-translation}
\end{equation}

Require \(\IF_t>0\) for every \(t\in\mathcal T_n\); convexity then makes the
population minimizer unique.  For \(\rho_H\), all three objectives are strictly
convex and \(\IF_t>0\) whenever \(W_t>0\).  They need not be globally strongly
convex: their curvatures tend to zero in the tails.  The proof uses positive
curvature near \(S_t\), not a uniform curvature lower bound on \(\R\).

The target \(S_t\) minimizes the expected noisy loss.  It generally differs
from \(\BiasedSignal_t\).  No assumption that the first component equals the full signal
value \(X_t\) is needed.

The constants below separate three controls: \(C_1\) bounds the centered
score, \(C_2\) bounds the curvature fluctuation, and \(C_3\) bounds the third
derivative.  The scale \(D_t=\sqrt{F_t}\) expresses the first two bounds in
local curvature units.  For \(\xx\geq0\) and
\(z_{\xx}\eqdef\xx+\log(4m)\), set
\begin{align}
  C_1
  &\eqdef
    \max_{t\in\mathcal T_n}
    \frac{L_1}{\DFN_t}
    \left\{Q_t\sqrt{2z_{\xx}}+\frac{2k_tz_{\xx}}{3}\right\},
  \label{eq:C1}\\
  C_2
  &\eqdef
    \max_{t\in\mathcal T_n}
    \frac{L_2}{\DFN_t}
    \left\{Q_t\sqrt{2z_{\xx}}+\frac{k_tz_{\xx}}{3}\right\},
  \label{eq:C2}\\
  C_3
  &\eqdef
    \max_{t\in\mathcal T_n}\frac{L_3W_t}{\DFN_t^2},
  &
  \DFN_{\min}
  &\eqdef\min_{t\in\mathcal T_n}\DFN_t,
  &
  \delta
  &\eqdef\frac{C_2}{\DFN_{\min}}.
  \label{eq:C3-delta}
\end{align}

The parameter \(\delta\) bounds the relative curvature error.  The two
conditions below keep the empirical curvature positive near \(S_t\) and
make the scalar Taylor expansion valid at the empirical minimizer.

\begin{proposition}[Sampling error expansion]\label{prop:main}
Suppose that the conditions above hold and that
\begin{equation}
  \delta<1,
  \qquad
  \frac{C_1C_3}{\DFN_{\min}(1-\delta)^2}<\frac49.
  \label{eq:smallness}
\end{equation}
Then there is a random set \(\Omega(\xx)\) with
\(\Pp\{\Omega(\xx)\}\geq1-e^{-\xx}\) such that
\begin{equation}
  \sup_{t\in\mathcal T_n}
  \left|
    \DFN_t\bigl(
      \IMFt_t-\IMFs_t-\IF_t^{-1}\nabla_t
    \bigr)
  \right|
  \leq
  \frac{1}{\DFN_{\min}}
  \left\{
    \frac{9C_1^2C_3}{8(1-\delta)^3}
    +\frac{C_1C_2}{1-\delta}
  \right\}.
  \label{eq:proposition-bound}
\end{equation}
On the same event, every empirical minimizer is unique and, for every \(t\),
\begin{equation}
  |\widetilde S_t-S_t|
    \leq\frac{3C_1}{2D_t(1-\delta)}.
  \label{eq:raw-localization}
\end{equation}
Define \(R_t=\widetilde S_t-S_t-F_t^{-1}\nabla_t\).  The
oracle-corrected estimator
\begin{equation}
  \IMFt_t^{\mathrm{corr}}
  \eqdef \IMFt_t-\IF_t^{-1}\nabla_t
  \label{eq:corrected-estimator}
\end{equation}
satisfies \(\widetilde S_t^{\mathrm{corr}}-S_t=R_t\), with the remainder
bound \eqref{eq:proposition-bound}.
\end{proposition}

\begin{remark}[Sign of the correction]
With \(\nabla_t\) defined by \eqref{eq:nabla}, the minus sign in
\eqref{eq:proposition-bound} is forced by differentiation:
\(f_\circ'(\IMFs_t)=-\nabla_t\).  Thus
\(\IMFt_t-\IMFs_t\) has first-order term \(\IF_t^{-1}\nabla_t\), and one
subtracts that term to correct \(\IMFt_t\).  This removes the leading
sampling fluctuation; it does not remove \(\Delta_t\).
\end{remark}

\subsection{Score and curvature bounds}\label{sec:derivatives}

The next lemma supplies the probabilistic input to
Proposition~\ref{prop:main}.  Its score and curvature bounds use Bernstein's
inequality and a union bound over the finite evaluation grid.  The third
derivative bounds are deterministic.

Fix \(t\in\mathcal T_n\).  All quantities in this subsection may depend on
\(t\).  For the proof, abbreviate the population and empirical objectives as
\begin{equation}
  \fs(z)\eqdef\sum_u\rhohum(\Signal_u-z)w_{u,t},
  \qquad
  \fp(z)\eqdef\sum_u\rhohu(\Data_u-z)w_{u,t},
  \label{eq:objectives}
\end{equation}
and
\begin{equation}
  \xs\eqdef\arginf_z\fs(z)=\IMFs_t,
  \qquad
  \xp\eqdef\arginf_z\fp(z)=\IMFt_t.
  \label{eq:minimizers}
\end{equation}
Further, put
\begin{equation}
  \IFP\eqdef\fp''(\xs),
  \qquad
  \IF\eqdef\fs''(\xs)=\DFN^2.
  \label{eq:F-Fcirc}
\end{equation}
For \(r>0\), let
\(U_t(r)\eqdef\{z\in\R:|\DFN(z-\xs)|\leq r\}\).

\begin{lemma}\label{lem:local}
There is an event \(\Omega(\xx)\) of probability at least \(1-e^{-\xx}\)
on which, simultaneously for every \(t\in\mathcal T_n\),
\begin{enumerate}
\renewcommand{\labelenumi}{(\roman{enumi})}
\item
\begin{align}
  \frac{1}{\DFN}
  |\fs'(\xs)-\fp'(\xs)|
  &=\frac{1}{\DFN}
    \left|\sum_u
      \left\{\rhohu'(\Data_u-\xs)
      -\E\rhohu'(\Data_u-\xs)\right\}w_{u,t}
    \right|
  \leq C_1,
  \label{eq:lemma-score}\\
  \frac{1}{\DFN}
  |\fs''(\xs)-\fp''(\xs)|
  &=\frac{1}{\DFN}
    \left|\sum_u
      \left\{\rhohu''(\Data_u-\xs)
      -\E\rhohu''(\Data_u-\xs)\right\}w_{u,t}
    \right|
  \leq C_2.
  \label{eq:lemma-hessian}
\end{align}

\item
\begin{equation}
  \sup_{z\in\R}
  \frac{1}{\DFN^2}
  \left|\sum_u\rhohu'''(\Data_u-z)w_{u,t}\right|
  \leq C_3.
  \label{eq:lemma-third}
\end{equation}

\item For every \(r>0\),
\begin{equation}
  \sup_{z\in U_t(r)}
  \frac{1}{\DFN^2}
  \left|\sum_u
    \left\{
      \rhohum''(\Signal_u-z)
      -\rhohum''(\Signal_u-\xs)
    \right\}w_{u,t}
  \right|
  \leq \frac{C_3r}{\DFN}.
  \label{eq:lemma-population-hessian}
\end{equation}
\end{enumerate}
\end{lemma}


The proof in Appendix~\ref{sec:scalar-proof} uses only the bounds in
Lemma~\ref{lem:local}.  It also applies with other constants satisfying those
bounds on a specified event.  Gaussian noise permits sharper variance bounds.


\subsection{Gaussian noise and normalized weights}\label{sec:gaussian}

Multiplying all weights in a window by a positive constant changes neither
minimizer nor \(F_t^{-1}\nabla_t\).  Define
\begin{equation}
  a_{u,t}\eqdef\frac{w_{u,t}}{W_t},\qquad
  N\eqdef\min_t\frac{1}{\sum_u a_{u,t}^2},\qquad
  \alpha\eqdef\max_{t,u}a_{u,t}.
  \label{eq:normalized-weights}
\end{equation}
Here \(N\) is the smallest effective sample size across windows.  The population
curvature and score for the normalized objective are
\(\bar F_t=F_t/W_t\) and \(\bar\nabla_t=\nabla_t/W_t\).
Their ratio is still \(F_t^{-1}\nabla_t\).

Assume now \(\varepsilon_u\sim N(0,\sigma^2)\) independently, with
\(\sigma\geq0\), and \(\rho=\rho_H\).  Put
\begin{equation}
  c\eqdef\sqrt{2/\pi},\qquad B\eqdef\sqrt{H^2+\sigma^2}.
  \label{eq:gaussian-scale}
\end{equation}
Independence of the Gaussian variables in the two expectations gives the exact
identity
\begin{equation}
  \phi(r)=\E_\varepsilon\rho_H(r+\varepsilon)=\rho_B(r).
  \label{eq:gaussian-convolution}
\end{equation}
Thus \(\BiasedSignal_t\) fits the clean signal at scale \(H\), whereas \(S_t\) fits it
at scale \(B\).  In particular,
\[
  \sum_u a_{u,t}\operatorname{erf}\!\left(\frac{X_u-\BiasedSignal_t}{\sqrt2H}\right)=0,
\]
and the population target and normalized curvature satisfy
\begin{equation}
  \sum_u a_{u,t}\operatorname{erf}\!\left(\frac{X_u-S_t}{\sqrt2B}\right)=0,
  \qquad
  \bar F_t=\frac cB\sum_u a_{u,t}e^{-(X_u-S_t)^2/(2B^2)}.
  \label{eq:gaussian-target}
\end{equation}
The weighted score is positive below all active signal values and negative
above them.  Its unique zero \(S_t\) therefore lies between the smallest and
largest \(X_u\) having positive weight.  If
\begin{equation}
  A\geq\max_t\max_{u,v:\,a_{u,t}a_{v,t}>0}|X_u-X_v|,
  \qquad
  \lambda\eqdef\frac cB e^{-A^2/(2B^2)},
  \label{eq:curvature-lower-bound}
\end{equation}
then \(\bar F_t\geq\lambda>0\) at every evaluation point.  More generally,
any known positive lower bound \(\lambda\leq\min_t\bar F_t\) can be used below.

The bounds \(s_1\) and \(s_2\) below control the normalized score and
curvature fluctuations.  Dividing \(s_2\) by the curvature lower bound
\(\lambda\) gives the relative perturbation \(d\).  For
\(L_2=c/H\) and \(L_3=c/H^2\), define
\begin{align}
  v_1&\eqdef\min\{1,\sigma L_2\},&
  v_2&\eqdef\min\{L_2,\sigma L_3\},\notag\\
  s_1&\eqdef v_1\sqrt{\frac{2z_x}{N}}+\frac{2\alpha z_x}{3},&
  s_2&\eqdef v_2\sqrt{\frac{2z_x}{N}}+\frac{L_2\alpha z_x}{3},\notag\\
  d&\eqdef\frac{s_2}{\lambda},&
  z_x&=x+\log(4m).
  \label{eq:gaussian-tail-constants}
\end{align}

\begin{corollary}[Explicit Gaussian bounds]\label{cor:gaussian}
Under the Gaussian model and the deterministic design conditions above, suppose
\begin{equation}
  d<1,\qquad
  \frac{s_1L_3}{\lambda^2(1-d)^2}<\frac49.
  \label{eq:gaussian-smallness}
\end{equation}
With probability at least \(1-e^{-x}\), simultaneously for all \(t\),
\begin{align}
  |\widetilde S_t-S_t|
    &\leq\mathcal E_x\eqdef\frac{3s_1}{2\lambda(1-d)},
      \label{eq:gaussian-raw}\\
  |F_t^{-1}\nabla_t|&\leq\frac{s_1}{\lambda},
      \label{eq:gaussian-leading}\\
  |R_t|&\leq\mathcal R_x
     \eqdef\frac{9L_3s_1^2}{8\lambda^3(1-d)^3}
            +\frac{s_1s_2}{\lambda^2(1-d)},
      \label{eq:gaussian-remainder}\\
  R_t&\eqdef\widetilde S_t-S_t-F_t^{-1}\nabla_t.
      \notag
\end{align}
In particular,
\begin{equation}
  |\widetilde S_t^{\mathrm{corr}}-S_t|\leq\mathcal R_x.
  \label{eq:gaussian-corrected}
\end{equation}
The constants \(\mathcal E_x\) and \(\mathcal R_x\) are specified by the
weights, \(x\), \(H\), \(\sigma\), and a curvature lower bound such as
\eqref{eq:curvature-lower-bound}.  These bounds concern the sampling error.
The target shift will enter only when we change the reference to
\(\BiasedSignal_t\) in Section~\ref{sec:biased}.
\end{corollary}


\begin{remark}[Dependence on the noise and loss scales]
The noise scale enters through \(B=\sqrt{H^2+\sigma^2}\), the lower bound
\(\lambda\), and the variance bounds \(v_1,v_2\).  These are sufficient,
conservative constants.  At \(\sigma=0\), the empirical and population
objectives coincide with the clean objective, so both errors and \(R_t\) vanish;
the Bernstein linear terms need not vanish.  If \(H\) varies with sample size,
the explicit conditions must be checked again: \(L_2\) and \(L_3\) grow as
\(H^{-1}\) and \(H^{-2}\) when \(H\) decreases.
\end{remark}

\subsection{Bandwidth and sample-size conditions}\label{sec:bandwidth}

Suppose that, for a known constant \(C_K>0\),
\begin{equation}
  N\geq\frac{nh}{C_K},\qquad \alpha\leq\frac{C_K}{nh}.
  \label{eq:kernel-mass}
\end{equation}
These assumptions say that the effective sample size is at least a constant
times \(nh\), and no individual observation has disproportionate weight.
They can be checked directly from the normalized weights.

For an equally spaced design, the calculation is clearest in terms of the
actual grid spacing \(\Delta\).  Use untruncated interior windows, or a
periodic grid with each observation counted once, and take \(0<h<1/2\).
Let \(w_{u,t}=K((t-u)/h)\), using local offsets in the periodic case, with
\(K\) supported on \([-1,1]\), \(0\leq K\leq K_{\max}\), and
\(K(v)\geq k_0>0\) on \(|v|\leq\eta\), where \(0<\eta<1\).
If \(\eta h\geq\Delta\), the central interval has length \(2\eta h\) and
contains at least
\[
  \left\lfloor\frac{2\eta h}{\Delta}\right\rfloor
  \geq\frac{\eta h}{\Delta}
\]
grid points.  Consequently,
\[
  W_t\geq\frac{k_0\eta h}{\Delta},\qquad
  \max_u a_{u,t}\leq\frac{K_{\max}\Delta}{k_0\eta h},\qquad
  \sum_u a_{u,t}^2\leq\max_u a_{u,t}.
\]
For the usual unit-interval grids, \(\Delta=1/n\) when the right endpoint is
excluded, and \(\Delta=1/(n-1)\) when both endpoints are included.  For
\(n\geq2\), both satisfy \(\Delta\leq2/n\).  Thus the single choice
\[
  C_K=\frac{2K_{\max}}{k_0\eta}
\]
establishes \eqref{eq:kernel-mass} for either convention, provided
\(\eta h\geq\Delta\).  The condition \(nh\geq2/\eta\) is sufficient.
If the spacing is exactly \(1/n\), the sharper constant
\(C_K=K_{\max}/(k_0\eta)\) and condition \(nh\geq1/\eta\) suffice.
A periodic grid has \(n\) distinct locations per period; it does not count
both endpoints as separate locations.  There are at most \(2h/\Delta+1\)
active observations, which also bounds the effective sample size from above.
Both grid conventions therefore give effective sample size of order \(nh\).

Set
\begin{equation}
  q_n\eqdef\sqrt{\frac{C_Kz_x}{nh}},\qquad
  A_1\eqdef\sqrt2v_1+\frac23,\qquad
  A_2\eqdef\sqrt2v_2+\frac{L_2}{3}.
  \label{eq:rate-constants}
\end{equation}

\begin{corollary}[A sufficient local sample size]\label{cor:sample-size}
Under the Gaussian model, the curvature lower bound, and
\eqref{eq:kernel-mass}, the following condition implies
\eqref{eq:gaussian-smallness}:
\begin{equation}
  nh>C_Kz_x\max\left\{
       1,\frac{4A_2^2}{\lambda^2},
       \frac{81A_1^2L_3^2}{\lambda^4}\right\}.
  \label{eq:sample-size-condition}
\end{equation}
With probability at least \(1-e^{-x}\),
\begin{align}
  \max_t|\widetilde S_t-S_t|
    &\leq\frac{3A_1}{\lambda}q_n,
      \label{eq:raw-rate}\\
  \max_t|R_t|
    &\leq\left\{\frac{9L_3A_1^2}{\lambda^3}
                 +\frac{2A_1A_2}{\lambda^2}\right\}q_n^2.
      \label{eq:remainder-rate}
\end{align}
\end{corollary}


For fixed \(H>0\), fixed \(\sigma\), bounded local signal oscillation, and
fixed kernel constants, \(\lambda\) stays bounded away from zero.  The two
bounds then have orders \(\sqrt{(x+\log(4m))/(nh)}\) and
\((x+\log(4m))/(nh)\), respectively.  The smallness conditions hold eventually
if \((x+\log(4m))/(nh)\to0\).  These are upper bounds on sampling error;
they do not assert equality of the remainder and the squared realized error.
The biased signal error also contains \(\Delta_t\), which need not decrease
when more observations are added at a fixed filtering scale.

\subsection{Target shift: comparing the deterministic targets}\label{sec:shift}

The preceding bounds isolate the sampling fluctuation.  A bound on
\(\widetilde S_t-\BiasedSignal_t\) also needs a bound on \(S_t-\BiasedSignal_t\).  Gaussian averaging
changes the loss scale from \(H\) to \(B\), so this shift can persist even
when the sampling fluctuation becomes small.

Let the active signal range and its curvature bound be
\begin{equation}
  A_t\eqdef\max_{u:a_{u,t}>0}X_u-\min_{u:a_{u,t}>0}X_u,
  \qquad
  \kappa_t\eqdef\frac cB e^{-A_t^2/(2B^2)}.
  \label{eq:shift-constants}
\end{equation}
Both \(S_t\) and \(\BiasedSignal_t\) lie in this signal range.  Throughout the interval
between them, the normalized population curvature is at least \(\kappa_t\).
This interval bound is stronger than knowing the curvature only at \(S_t\).

\begin{proposition}[Deterministic target shift]\label{prop:shift}
Under the Gaussian model, put
\begin{equation}
  \beta_t\eqdef\min\!\left\{A_t,
                  \frac{L_3\sigma^2}{2\kappa_t}\right\},
  \qquad \beta\eqdef\max_t\beta_t.
  \label{eq:beta}
\end{equation}
Then, deterministically,
\begin{equation}
  |S_t-\BiasedSignal_t|\leq\beta_t\qquad\text{for every }t.
  \label{eq:target-shift-bound}
\end{equation}
One may take \(\lambda=\min_t\kappa_t>0\) in
Corollary~\ref{cor:gaussian}.
\end{proposition}

For fixed signal, weights, and \(H\), the shift bound is \(O(\sigma^2)\)
as \(\sigma\to0\).  It is a sufficient upper bound and can be conservative.
At \(\sigma=0\), both errors vanish exactly.  If the weighted signal values
are symmetric about a number \(\mu_t\), both target scores vanish there,
so \(\BiasedSignal_t=S_t=\mu_t\) even when \(\sigma>0\).  This includes constant
signals.  Symmetric observation locations alone do not imply symmetry of
signal values for a curved signal.

\begin{remark}[An asymmetric window]
Take equally weighted signal values \((0,0,3)\), \(H=0.5\), and
\(\sigma=1\).  The target at scale \(b\) is the unique root on \([0,3]\) of
\[
  2\operatorname{erf}\!\left(\frac{-z}{\sqrt2b}\right)
  +\operatorname{erf}\!\left(\frac{3-z}{\sqrt2b}\right)=0.
\]
Bisection at \(b=H\) and \(b=B=\sqrt{1.25}\) gives
\[
  \BiasedSignal_t\approx0.337245,\qquad S_t\approx0.718197,
  \qquad\Delta_t\approx0.380952.
\]
Repeating each signal value \(r\) times, with weights \(1/(3r)\) and
independent observation noise, preserves both targets and gives effective
sample size \(3r\).  As \(r\to\infty\), the sampling error tends to zero
in probability, while the biased signal error tends to \(\Delta_t\).
\end{remark}

\subsection{Biased signal error: adding the target shift}\label{sec:biased}

The sampling error and target shift now have separate bounds.  The identity
\eqref{eq:target-translation} combines them:
\begin{equation}
  \boxed{\widetilde S_t-\BiasedSignal_t
    =\underbrace{F_t^{-1}\nabla_t+R_t}_{\text{sampling error}}
     +\underbrace{\Delta_t}_{\text{target shift}}.}
  \label{eq:biased-expansion}
\end{equation}
The same shift appears after oracle correction:
\(\widetilde S_t^{\mathrm{corr}}-\BiasedSignal_t=R_t+\Delta_t\).

\begin{corollary}[Combined error bounds]\label{cor:biased}
Under the Gaussian model, if the smallness conditions of Corollary~\ref{cor:gaussian} hold, then with
probability at least \(1-e^{-x}\), simultaneously for all \(t\),
\begin{align}
  |\widetilde S_t-S_t|&\leq\mathcal E_x,
    \label{eq:final-population}\\
  |\widetilde S_t-\BiasedSignal_t|&\leq\beta_t+\mathcal E_x,
    \label{eq:final-biased}\\
  |\widetilde S_t^{\mathrm{corr}}-S_t|&\leq\mathcal R_x,
    \label{eq:final-corrected-population}\\
  |\widetilde S_t^{\mathrm{corr}}-\BiasedSignal_t|&\leq\beta_t+\mathcal R_x.
    \label{eq:final-corrected-biased}
\end{align}
Replacing \(\beta_t\) by \(\beta\) gives uniform bounds over the evaluation
grid.  The shift bound is deterministic, so adding it does not change the
failure probability.
\end{corollary}

Under the additional conditions of Corollary~\ref{cor:sample-size}, the
same argument gives, with probability at least \(1-e^{-x}\),
\begin{align}
  \max_t|\widetilde S_t-\BiasedSignal_t|
    &\leq\beta+\frac{3A_1}{\lambda}q_n,
      \label{eq:biased-rate}\\
  \max_t|\widetilde S_t^{\mathrm{corr}}-\BiasedSignal_t|
    &\leq\beta+\left\{\frac{9L_3A_1^2}{\lambda^3}
                 +\frac{2A_1A_2}{\lambda^2}\right\}q_n^2.
      \label{eq:corrected-biased-rate}
\end{align}
Thus the sampling contribution decreases at the stated rates, while the
deterministic target shift remains.

For any realization, the reverse triangle inequality also gives
\begin{equation}
 \left|\max_t|\widetilde S_t-\BiasedSignal_t|-\max_t|\Delta_t|\right|
 \leq\max_t|\widetilde S_t-S_t|.
 \label{eq:error-floor}
\end{equation}
Thus a nonzero shift remains as sampling error decreases.  The upper bound
\(\beta\) itself is not a lower bound on the error.

\subsection{Scope of the conclusions}

All uniform probability statements concern the finite evaluation grid.  A
supremum over the continuum \([0,1]\) requires additional regularity and
control of the family of weighted processes.  Adaptive choices of bandwidth
or loss scale also require simultaneous control over those choices.

The oracle correction uses the unknown population target, noise law, and
curvature.  It is an expansion tool, not an estimator computable from the
observations alone.  Its remainder bound around \(S_t\) becomes a bound
around \(\BiasedSignal_t\) only after adding the target shift.  Neither signed error is
guaranteed to have the smaller absolute value on every realization.

In synthetic experiments, compute \(\BiasedSignal\) by applying the filter to the clean
signal at scale \(H\), and compute \(S\) at scale \(B\) under Gaussian
noise.  Evaluate both errors on the same noise realization and report
\(S-\BiasedSignal\).  Comparing with \(S\) tests the stochastic expansion; comparing
with \(\BiasedSignal\) measures how noise changes the biased signal.  Neither
comparison alone establishes recovery of latent generative components.

The general proposition allows independent non-Gaussian errors with a finite
first moment.  The explicit Gaussian formulas for the population target and
its shift do not cover arbitrary non-Gaussian noise.  Later IRMF components
use estimated residuals that share observations, so their dependent errors
require a separate argument.

Throughout, \(F_t\) is evaluated at \(S_t\).  It cannot be replaced by a
curvature evaluated at \(\BiasedSignal_t\) without bounding their difference.


\appendix
\section{Proofs}\label{sec:proofs}

The argument has two stages.  Lemma~\ref{lem:local} controls the random score
and curvature at the population target.  Conditional on that event, the proof
of Proposition~\ref{prop:main} is deterministic: localize the empirical
minimizer, apply scalar Taylor expansion, and replace empirical curvature by
population curvature.  The remaining proofs specialize these bounds and
control the target shift.

\subsection{Score and curvature concentration}

\begin{proof}[Proof of Lemma~\ref{lem:local}]
Differentiating \eqref{eq:objectives} gives
\begin{align*}
  \fs'(z)&=-\sum_u\rhohum'(\Signal_u-z)w_{u,t},
  &
  \fp'(z)&=-\sum_u\rhohu'(\Data_u-z)w_{u,t},\\
  \fs''(z)&=\sum_u\rhohum''(\Signal_u-z)w_{u,t},
  &
  \fp''(z)&=\sum_u\rhohu''(\Data_u-z)w_{u,t}.
\end{align*}
Since differentiation under the expectation is justified by
\eqref{eq:derivative-bounds},
\(\rhohum^{(j)}(a)=\E\rhohu^{(j)}(a+\varepsilon)\) for
\(j=1,2,3\).  The identities in part~(i) follow.

For the first identity, set
\[
  Z_{u,t}\eqdef w_{u,t}
  \left\{\rhohu'(\Data_u-\xs)
  -\E\rhohu'(\Data_u-\xs)\right\}.
\]
These variables are independent and centered.  Moreover,
\[
  |Z_{u,t}|\leq2L_1k_t,
  \qquad
  \sum_u\E Z_{u,t}^2\leq L_1^2Q_t^2.
\]
Bernstein's inequality therefore yields, for every \(z>0\),
\begin{equation}
  \Pp\left(
    \left|\sum_uZ_{u,t}\right|>
    L_1\left\{Q_t\sqrt{2z}+\frac{2k_tz}{3}\right\}
  \right)
  \leq2e^{-z}.
  \label{eq:bernstein-score}
\end{equation}
For the second identity, use
\[
  \widetilde Z_{u,t}\eqdef w_{u,t}
  \left\{\rhohu''(\Data_u-\xs)
  -\E\rhohu''(\Data_u-\xs)\right\}.
\]
Convexity and \eqref{eq:derivative-bounds} give
\[
  |\widetilde Z_{u,t}|\leq L_2k_t,
  \qquad
  \sum_u\E\widetilde Z_{u,t}^2\leq L_2^2Q_t^2.
\]
Hence
\begin{equation}
  \Pp\left(
    \left|\sum_u\widetilde Z_{u,t}\right|>
    L_2\left\{Q_t\sqrt{2z}+\frac{k_tz}{3}\right\}
  \right)
  \leq2e^{-z}.
  \label{eq:bernstein-hessian}
\end{equation}
Take \(z=z_{\xx}=\xx+\log(4m)\) in
\eqref{eq:bernstein-score} and \eqref{eq:bernstein-hessian}, and take the
union over the two inequalities and over \(t\in\mathcal T_n\).  The exceptional
probability is at most
\(4m e^{-z_{\xx}}=e^{-\xx}\).  Equations \eqref{eq:C1} and
\eqref{eq:C2} now prove part~(i).

Part~(ii) is deterministic.  By \eqref{eq:derivative-bounds} and
\eqref{eq:C3-delta},
\[
  \sup_{z\in\R}\frac{1}{\DFN^2}
  \left|\sum_u\rhohu'''(\Data_u-z)w_{u,t}\right|
  \leq\frac{L_3W_t}{\DFN^2}\leq C_3.
\]
Finally, \(|\rhohum'''|\leq L_3\).  The mean value theorem and
\(|z-\xs|\leq r/\DFN\) for \(z\in U_t(r)\) give
\begin{align*}
  &\frac{1}{\DFN^2}
  \left|\sum_u
    \left\{\rhohum''(\Signal_u-z)
    -\rhohum''(\Signal_u-\xs)\right\}w_{u,t}\right|\\
  &\hspace{35mm}\leq
    \frac{L_3W_t}{\DFN^2}|z-\xs|
    \leq\frac{C_3r}{\DFN},
\end{align*}
which proves part~(iii).
\end{proof}

\subsection{Scalar localization and the remainder}\label{sec:scalar-proof}

\begin{proof}[Proof of Proposition~\ref{prop:main}]
Fix \(t\in\mathcal T_n\).  Since \(\fs'(\xs)=0\),
\begin{equation}
  \fp'(\xs)
  =-\sum_u\left\{
      \rhohu'(\Data_u-\xs)-\E\rhohu'(\Data_u-\xs)
    \right\}w_{u,t}
  =-\nabla_t.
  \label{eq:score-sign}
\end{equation}
Lemma~\ref{lem:local}(i) also gives
\begin{equation}
  \frac{|\IFP-\IF|}{\IF}
  \leq\frac{C_2}{\DFN}
  \leq\delta.
  \label{eq:F-comparison}
\end{equation}
In particular, \(\IFP\geq(1-\delta)\IF>0\).  With
\begin{equation}
  b_{\circ}\eqdef
  \frac{\DFN}{\IFP}|\fp'(\xs)|,
  \label{eq:bcirc}
\end{equation}
the score bound and \eqref{eq:F-comparison} imply
\begin{equation}
  b_{\circ}\leq\frac{C_1}{1-\delta}.
  \label{eq:bcirc-bound}
\end{equation}

Use the scaled coordinate \(v=D(z-x^*)\), and put
\begin{equation}
  g(v)\eqdef f_\circ(x^*+v/D),\qquad
  a\eqdef g''(0)=F_\circ/D^2,\qquad
  \tau\eqdef C_3/D.
  \label{eq:scaled-objective}
\end{equation}
Then \(a\geq1-\delta\), \(|g'''(v)|\leq\tau\), and
\(b_\circ=|g'(0)|/a\).  Consequently,
\begin{equation}
  \frac{\tau b_\circ}{a}
  \leq\frac{C_1C_3}{D(1-\delta)^2}<\frac49.
  \label{eq:scaled-smallness}
\end{equation}
If \(b_\circ=0\), convexity and \(g''(0)>0\) imply that \(0\) is the unique
minimizer, and the claimed bounds are immediate.  Otherwise set
\(r=3b_\circ/2\).  Taylor expansion of \(g'\) at zero yields
\[
  |g'(v)-g'(0)-av|\leq\frac{\tau}{2}v^2.
\]
Since \(|g'(0)|=ab_\circ\),
\begin{align*}
  g'(r)&\geq-ab_\circ+ar-\frac{\tau r^2}{2}
       =\frac{ab_\circ}{2}-\frac{9\tau b_\circ^2}{8}>0,\\
  g'(-r)&\leq ab_\circ-ar+\frac{\tau r^2}{2}
       =-\frac{ab_\circ}{2}+\frac{9\tau b_\circ^2}{8}<0.
\end{align*}
Moreover, for \(|v|\leq r\),
\[
  g''(v)\geq a-\tau r>\frac a3>0.
\]
There is therefore a unique zero \(v_\circ\in(-r,r)\) of \(g'\).
Convexity on \(\R\) makes it the unique global minimizer: the derivative is
negative to the left of this interval and positive to its right.  This proves
\begin{equation}
  |v_\circ|=D|x_\circ-x^*|\leq\frac32b_\circ
  \leq\frac{3C_1}{2(1-\delta)}.
  \label{eq:scalar-localization}
\end{equation}
At the zero of \(g'\), the same Taylor expansion gives
\begin{equation}
  \left|a\left(v_\circ+\frac{g'(0)}a\right)\right|
  \leq\frac{\tau}{2}v_\circ^2
  \leq\frac{9\tau}{8}b_\circ^2.
  \label{eq:scalar-remainder}
\end{equation}
Equivalently, the left-hand side is
\(D^{-1}F_\circ|x_\circ-x^*+F_\circ^{-1}f_\circ'(x^*)|\).
Using \eqref{eq:F-comparison} and \eqref{eq:bcirc-bound}, we obtain
\begin{equation}
  \left|\xp-\xs+\IFP^{-1}\fp'(\xs)\right|
  \leq
  \frac{9C_1^2C_3}{8\DFN^2(1-\delta)^3}.
  \label{eq:Fcirc-remainder}
\end{equation}

It remains to replace \(\IFP\) by the population curvature \(\IF\).  By
\eqref{eq:F-comparison}, \eqref{eq:lemma-score}, and
\(\fs'(\xs)=0\),
\begin{align}
  \left|\IFP^{-1}\fp'(\xs)-\IF^{-1}\fp'(\xs)\right|
  &=\frac{|\IFP-\IF|}{\IFP\IF}|\fp'(\xs)|
  \notag\\
  &\leq\frac{C_1C_2}{\DFN^2(1-\delta)}.
  \label{eq:F-replacement}
\end{align}
The triangle inequality applied to \eqref{eq:Fcirc-remainder} and
\eqref{eq:F-replacement}, followed by \eqref{eq:score-sign}, proves
\begin{equation}
  \DFN\left|
    \xp-\xs-\IF^{-1}\nabla_t
  \right|
  \leq
  \frac{1}{\DFN}
  \left\{
    \frac{9C_1^2C_3}{8(1-\delta)^3}
    +\frac{C_1C_2}{1-\delta}
  \right\}.
  \label{eq:pointwise-final}
\end{equation}
Taking the supremum over \(t\in\mathcal T_n\) and using
\(\DFN_t\geq\DFN_{\min}\) gives \eqref{eq:proposition-bound}.
Subtracting \(F_t^{-1}\nabla_t\) gives
\(\widetilde S_t^{\mathrm{corr}}-S_t=R_t\).
\end{proof}

\begin{remark}[The remainder coefficient]
The coefficient \(9/8\) is \(\tfrac12(\tfrac32)^2\).  The factor \(1/2\)
comes from Taylor's theorem and \(3/2\) from the localization radius.
\end{remark}

\subsection{Gaussian constants}

\begin{proof}[Proof of Corollary~\ref{cor:gaussian}]
The identity \eqref{eq:gaussian-convolution} follows because
\(\varepsilon+HZ\sim N(0,B^2)\); the target and curvature formulas follow by
differentiation.  For a Lipschitz function \(G\) and an independent copy
\(\varepsilon'\) of \(\varepsilon\),
\[
  \operatorname{Var}\{G(X_u+\varepsilon)\}
  =\tfrac12\E\{G(X_u+\varepsilon)-G(X_u+\varepsilon')\}^2
  \leq\operatorname{Lip}(G)^2\sigma^2.
\]
Apply this with \(G(y)=\rho_H'(y-S_t)\) and
\(G(y)=\rho_H''(y-S_t)\).  Their Lipschitz constants are at most \(L_2\)
and \(L_3\).  Combining this estimate with the boundedness estimates from
Lemma~\ref{lem:local} gives variance bounds \(v_1^2\) and \(v_2^2\).
Bernstein's inequality and the same union bound therefore give an event of
probability at least \(1-e^{-x}\) on which
\[
  \max_t|\bar\nabla_t|\leq s_1,\qquad
  \max_t|\bar F_{\circ,t}-\bar F_t|\leq s_2,
  \qquad
  \bar F_{\circ,t}=\sum_u a_{u,t}\rho_H''(Y_u-S_t).
\]
For the normalized objectives, the deterministic proof of
Proposition~\ref{prop:main} applies with the bounds
\[
  C_1=\frac{s_1}{\sqrt\lambda},\qquad
  C_2=\frac{s_2}{\sqrt\lambda},\qquad
  C_3=\frac{L_3}{\lambda},\qquad
  D_{\min}\geq\sqrt\lambda,\qquad\delta\leq d.
\]
Condition \eqref{eq:gaussian-smallness} is sufficient for
\eqref{eq:smallness}.  Substitution in \eqref{eq:raw-localization}, using
\(D_t\geq\sqrt\lambda\), proves the raw bound.  Substitution in
\eqref{eq:pointwise-final}, followed by division by \(D_t\geq\sqrt\lambda\),
proves the remainder bound.  The score bound and
\(\bar F_t\geq\lambda\) give \eqref{eq:gaussian-leading}.
The identity \(\widetilde S_t^{\mathrm{corr}}-S_t=R_t\) gives
\eqref{eq:gaussian-corrected}.
\end{proof}

\subsection{Sufficient local sample size}

\begin{proof}[Proof of Corollary~\ref{cor:sample-size}]
The mass bounds imply
\[
  s_1\leq\sqrt2v_1q_n+\tfrac23q_n^2,\qquad
  s_2\leq\sqrt2v_2q_n+\tfrac{L_2}{3}q_n^2.
\]
Condition \eqref{eq:sample-size-condition} implies
\[
  q_n<\min\left\{1,\frac{\lambda}{2A_2},
                    \frac{\lambda^2}{9A_1L_3}\right\}.
\]
Therefore \(s_1\leq A_1q_n\), \(s_2\leq A_2q_n\), \(d<1/2\), and
\[
  \frac{s_1L_3}{\lambda^2(1-d)^2}
  \leq\frac{4A_1q_nL_3}{\lambda^2}<\frac49.
\]
Apply Corollary~\ref{cor:gaussian} and use
\((1-d)^{-1}\leq2\) and \((1-d)^{-3}\leq8\).
\end{proof}

\subsection{Target shift and biased signal error}

\begin{proof}[Proof of Proposition~\ref{prop:shift}]
Write
\[
  \psi_b(r)=\operatorname{erf}\!\left(\frac{r}{\sqrt2b}\right),
  \qquad G_{b,t}(z)=\sum_u a_{u,t}\psi_b(X_u-z).
\]
Both scores are strictly decreasing, positive below the active signal range,
and negative above it.  Their zeros \(\BiasedSignal_t\) and \(S_t\) therefore lie in
that range.  In particular, \(|S_t-\BiasedSignal_t|\leq A_t\).

Differentiating the Gaussian convolution identity gives
\(\psi_B(r)=\E\psi_H(r+\varepsilon)\).  Since
\(\sup_r|\psi_H''(r)|=\sup_r|\rho_H'''(r)|\leq L_3\), Taylor's theorem gives
\[
  |\psi_H(r+e)-\psi_H(r)-e\psi_H'(r)|\leq\frac{L_3}{2}e^2
  \qquad(r,e\in\R).
\]
Taking expectation with \(e=\varepsilon\), and using
\(\E\varepsilon=0\) and \(\E\varepsilon^2=\sigma^2\), yields
\[
  |\psi_B(r)-\psi_H(r)|\leq\frac{L_3\sigma^2}{2}.
\]
As \(G_{H,t}(\BiasedSignal_t)=0\) and the weights sum to one,
\begin{equation}
  |G_{B,t}(\BiasedSignal_t)|
    =\left|\sum_u a_{u,t}
      \{\psi_B(X_u-\BiasedSignal_t)-\psi_H(X_u-\BiasedSignal_t)\}\right|
    \leq\frac{L_3\sigma^2}{2}.
  \label{eq:shift-score}
\end{equation}

For every \(z\) between \(\BiasedSignal_t\) and \(S_t\), active residuals satisfy
\(|X_u-z|\leq A_t\).  Hence
\begin{equation}
  -G_{B,t}'(z)=\frac cB\sum_u a_{u,t}e^{-(X_u-z)^2/(2B^2)}
             \geq\kappa_t.
  \label{eq:shift-curvature}
\end{equation}
Since \(G_{B,t}(S_t)=0\), the mean value theorem gives
\[
  |S_t-\BiasedSignal_t|\leq\frac{|G_{B,t}(\BiasedSignal_t)|}{\kappa_t}
            \leq\frac{L_3\sigma^2}{2\kappa_t}.
\]
Combining this with the range bound proves \eqref{eq:target-shift-bound}.
The same curvature bound at \(S_t\) gives
\(\bar F_t\geq\kappa_t\), so \(\lambda=\min_t\kappa_t\) is a valid
choice in Corollary~\ref{cor:gaussian}.

\end{proof}

\begin{proof}[Proof of Corollary~\ref{cor:biased} and the rate bounds]
On the event of Corollary~\ref{cor:gaussian}, add the deterministic bound
\(|\Delta_t|\leq\beta_t\) to \eqref{eq:gaussian-raw} and to the corrected
remainder bound.  Identity \eqref{eq:biased-expansion} and the triangle
inequality give \eqref{eq:final-population}--\eqref{eq:final-corrected-biased}
without changing the exceptional probability.  Adding \(\beta\) to the
bounds in Corollary~\ref{cor:sample-size} gives
\eqref{eq:biased-rate} and \eqref{eq:corrected-biased-rate}.
\end{proof}

\end{document}
